Source : Insee, comptes de patrimoine annuels en base 2020 (tableaux 8.201 et 8.210, lignes « terrains supportant des bâtiments » et « logements », valeurs nettes de fin d'année), PIB à prix courants ; calculs de l'auteur.\par Lecture : fin 2024, les terrains sous les bâtiments valent 229 \% du PIB pour l'ensemble de l'économie ; pour les ménages, le terrain représente 45 \% de la valeur cumulée terrain + logement. La valeur du terrain est un résidu : valeur de marché de l'immobilier moins valeur de remplacement du bâti, de sorte que toute plus-value immobilière lui est imputée.
